\begin{figure}[htbp]
\centering
\resizebox{\linewidth}{!}{%
\begin{tikzpicture}[
  x=1cm, y=1cm,
  font=\sffamily,
  every node/.style={text=fluxTexto, align=center, inner sep=2pt},
  caixa/.style={rounded corners=6pt, line width=0.9pt},
  titulo/.style={font=\sffamily\large\bfseries},
  seta/.style={
    -{Latex[length=3mm, width=2.3mm]},
    line width=1pt, draw=fluxCinza,
    shorten <=1pt, shorten >=2pt
  }
]

% Blocos principais
\draw[caixa, draw=fluxAzul, fill=fluxAzul!8]
  (0,0) rectangle (4.4,3.6);
\draw[caixa, draw=fluxVerde, fill=fluxVerde!9]
  (5.75,0) rectangle (10.85,3.6);
\draw[caixa, draw=fluxRoxo, fill=fluxRoxo!8]
  (12.2,0) rectangle (19.2,3.6);
\draw[caixa, draw=fluxOcre, fill=fluxOcre!8]
  (0,-4.55) rectangle (6.5,-1.95);
\draw[caixa, draw=fluxCinza, fill=fluxCinza!7]
  (10.95,-4.55) rectangle (19.2,-1.95);

% Condicionamento ERA5
\node[titulo, text width=4cm] at (2.2,2.99)
  {Condicionamento\\ERA5};
\node[text=fluxAzul, font=\sffamily\LARGE] at (2.2,2.02)
  {$X_t$};
\node[font=\sffamily\small, text width=4cm] at (2.2,0.96)
  {Temperatura \textbullet{} Umidade\\
   Vento \textbullet{} Vapor na coluna\\
   12 canais meteorológicos};

% Geração condicional
\node[caixa, draw=fluxVerde, fill=white,
  inner sep=5pt, minimum height=8.5mm]
  (ruido) at (8.3,4.37) {Ruído de amostragem \(z\)};
\node[titulo, text width=4.7cm] at (8.3,2.95)
  {Geração condicional};
\node[font=\sffamily\LARGE\bfseries, text=fluxVerde]
  at (8.3,2.16) {CorrDiff};
\node[font=\sffamily\LARGE] at (8.3,1.07)
  {$p_{\theta}\bigl(Y_{s(t)} \mid X_t\bigr)$};

% Conjunto de realizações
\node[titulo, text width=6.6cm] at (15.7,3.0)
  {Conjunto de realizações};
\foreach \indice/\posicao in {1/13.68,2/15.70,3/17.72} {
  \node[caixa, draw=fluxRoxo, fill=white, text=fluxRoxo,
    minimum width=1.72cm, minimum height=1.24cm,
    font=\sffamily\LARGE]
    at (\posicao,1.75) {$\widehat{Y}_{s(t)}^{(\indice)}$};
}
\node[text width=6.6cm] at (15.7,0.53)
  {Amostras do modelo sob\\o mesmo condicionamento};

% Referência observada
\node[titulo, text width=6.1cm] at (3.25,-2.50)
  {Referência observada};
\node[font=\sffamily\large, text width=6.1cm] at (3.25,-3.25)
  {Radar \(Y_{s(t)}\) e máscara \(M_{s(t)}\)};
\node[text width=6.1cm] at (3.25,-4.04)
  {Pareamento: \(\lvert s(t)-t\rvert \leq 120\,\mathrm{s}\)};

% Treinamento e avaliação
\node[titulo, text width=7.8cm] at (15.075,-2.50)
  {Treinamento e avaliação};
\node[text width=7.8cm] at (15.075,-3.63)
  {Supervisão nas posições válidas\\[2pt]
   Intensidade \textbullet{} Estrutura espacial\\
   Qualidade probabilística};

% Setas e rótulos
\draw[seta] (4.4,1.8) -- (5.75,1.8);
\draw[seta, draw=fluxVerde] (ruido.south) -- (8.3,3.6);
\draw[seta] (10.85,1.8) -- (12.2,1.8);
\draw[seta] (15.7,0) -- (15.7,-1.95)
  node[midway, right=3pt] {Comparação};
\draw[seta] (6.5,-3.25) -- (10.95,-3.25);
\draw[seta, dashed, draw=fluxOcre] (6.5,-2.15) -- (7.38,0);
\node[anchor=west, align=left, text=fluxOcre] at (7.75,-1.28)
  {Supervisão\\no treinamento};

\end{tikzpicture}%
}
    \caption{Geração contemporânea condicionada ERA5–radar}
    \label{geracao-condicional}
\end{figure}